% Preámbulo del short paper en IEEEtran (conference, dos columnas).

% pdflatex no trae estos caracteres del texto: el signo menos (U+2212), gamma y delta.
\ifPDFTeX
  \DeclareUnicodeCharacter{2212}{\textminus}
  \DeclareUnicodeCharacter{03B3}{\ensuremath{\gamma}}
  \DeclareUnicodeCharacter{0394}{\ensuremath{\Delta}}
\fi
% Cortes de línea en cualquier punto de las URL largas de la bibliografía.
\usepackage{xurl}
\usepackage{booktabs,array}

% Títulos de figuras y tablas al estilo del IEEE: «Figura 1.» en letra chica; «TABLA I» centrado, con el título
% en versalitas debajo. Quarto carga el paquete caption para sus referencias cruzadas.
\usepackage{caption}
\DeclareCaptionLabelFormat{mayusculas}{\MakeUppercase{#1}~#2}
\captionsetup{font=footnotesize,labelsep=period}
\captionsetup[table]{labelformat=mayusculas,labelsep=newline,justification=centering,textfont=sc}

% El resumen del paper es el «Resumen ejecutivo» que pide el enunciado (babel lo llamaría «Resumen»).
\AtBeginDocument{\renewcommand{\abstractname}{Resumen ejecutivo}}

% Las tablas anchas van a lo ancho de la página (table*), que LaTeX solo coloca arriba de una página; las
% figuras, a una columna, cerca de donde se citan.
\setcounter{dbltopnumber}{2}
\renewcommand{\dbltopfraction}{0.85}
\renewcommand{\dblfloatpagefraction}{0.7}
\renewcommand{\topfraction}{0.85}
\renewcommand{\bottomfraction}{0.6}
\renewcommand{\textfraction}{0.1}
\renewcommand{\floatpagefraction}{0.75}
\setcounter{totalnumber}{4}
